\documentclass[a4paper]{jsarticle}
\usepackage{gaiyo2}
\usepackage{algorithmic}
\usepackage{amscd}
\usepackage{amsfonts}
\usepackage{amsmath}
\usepackage[psamsfonts]{amssymb}
\usepackage{amsthm}
\usepackage{ascmac}
\usepackage{bm}
\usepackage{color}
\usepackage{enumerate}
\usepackage{fancybox}
\usepackage[stable]{footmisc}
\usepackage{graphicx}
\usepackage{listings}
\usepackage{mathrsfs}
\usepackage{mathtools}
\usepackage{otf}
\usepackage{pifont}
\usepackage{proof}
\usepackage{subfigure}
\usepackage{tikz}
\usepackage{verbatim}
\usepackage[all]{xy}


\usetikzlibrary{cd}
\usetikzlibrary{arrows.meta}


% ================================
% パッケージを追加する場合のスペース 
\usepackage[dvipdfmx]{hyperref}
\usepackage{xcolor}
\definecolor{darkgreen}{rgb}{0,0.45,0} 
\definecolor{darkred}{rgb}{0.75,0,0}
\definecolor{darkblue}{rgb}{0,0,0.6} 
\hypersetup{
    colorlinks=true,
    %citecolor=darkgreen,
    %linkcolor=darkred,
    %urlcolor=darkblue,
    citecolor=black,
    linkcolor=black,
    urlcolor=black,
}
\usepackage{pxjahyper}

%\usepackage[mathscr]{euscript} % mathscr のフォントを変更

%\usepackage{mmasym}

%=================================


% --------------------------
% theoremstyle
% --------------------------
%\theoremstyle{definition}

% --------------------------
% newtheoem
% --------------------------

% 日本語で定理, 命題, 証明などを番号付きで用いるためのコマンドです. 
% If you want to use theorem environment in Japanece, 
% you can use these code. 
% Attention!
% All theorem enivironment numbers depend on 
% only section numbers.
\newtheorem{Axiom}{公理}[section]
\newtheorem{Definition}[Axiom]{定義}
\newtheorem{Theorem}[Axiom]{定理}
\newtheorem{Proposition}[Axiom]{命題}
\newtheorem{Lemma}[Axiom]{補題}
\newtheorem{Corollary}[Axiom]{系}
\newtheorem{Example}[Axiom]{例}
\newtheorem{Claim}[Axiom]{主張}
\newtheorem{Property}[Axiom]{性質}
\newtheorem{Attention}[Axiom]{注意}
\newtheorem{Question}[Axiom]{問}
\newtheorem{Problem}[Axiom]{問題}
\newtheorem{Consideration}[Axiom]{考察}
\newtheorem{Alert}[Axiom]{警告}
\newtheorem{Fact}[Axiom]{事実}


% 日本語で定理, 命題, 証明などを番号なしで用いるためのコマンドです. 
% If you want to use theorem environment with no number in Japanese, You can use these code.
\newtheorem*{Axiom*}{公理}
\newtheorem*{Definition*}{定義}
\newtheorem*{Theorem*}{定理}
\newtheorem*{Proposition*}{命題}
\newtheorem*{Lemma*}{補題}
\newtheorem*{Example*}{例}
\newtheorem*{Corollary*}{系}
\newtheorem*{Claim*}{主張}
\newtheorem*{Property*}{性質}
\newtheorem*{Attention*}{注意}
\newtheorem*{Question*}{問}
\newtheorem*{Problem*}{問題}
\newtheorem*{Consideration*}{考察}
\newtheorem*{Alert*}{警告}
\newtheorem{Fact*}{事実}


% 英語で定理, 命題, 証明などを番号付きで用いるためのコマンドです. 
%　If you want to use theorem environment in English, You can use these code.
%all theorem enivironment number depend on only section number.
\newtheorem{Axiom+}{Axiom}[section]
\newtheorem{Definition+}[Axiom+]{Definition}
\newtheorem{Theorem+}[Axiom+]{Theorem}
\newtheorem{Proposition+}[Axiom+]{Proposition}
\newtheorem{Lemma+}[Axiom+]{Lemma}
\newtheorem{Corollary+}[Axiom+]{Corollary}
\newtheorem{Claim+}[Axiom+]{Claim}
\newtheorem{Property+}[Axiom+]{Property}
\newtheorem{Attention+}[Axiom+]{Attention}
\newtheorem{Question+}[Axiom+]{Question}
\newtheorem{Problem+}[Axiom+]{Problem}
\newtheorem{Consideration+}[Axiom+]{Consideration}
\newtheorem{Alert+}{Alert}
\newtheorem{Fact+}[Axiom+]{Fact}

\theoremstyle{definition}
\newtheorem{Example+}[Axiom+]{Example}
\newtheorem{Remark+}[Axiom+]{Remark}

% ----------------------------
% commmand
% ----------------------------
% 執筆に便利なコマンド集です. 
% コマンドを追加する場合は下のスペースへ. 

% 集合の記号 (黒板文字)
\newcommand{\NN}{\mathbb{N}}
\newcommand{\ZZ}{\mathbb{Z}}
\newcommand{\QQ}{\mathbb{Q}}
\newcommand{\RR}{\mathbb{R}}
\newcommand{\CC}{\mathbb{C}}
\newcommand{\PP}{\mathbb{P}}
\newcommand{\KK}{\mathbb{K}}


% 集合の記号 (太文字)
\newcommand{\nn}{\mathbf{N}}
\newcommand{\zz}{\mathbf{Z}}
\newcommand{\qq}{\mathbf{Q}}
\newcommand{\rr}{\mathbf{R}}
\newcommand{\cc}{\mathbf{C}}
\newcommand{\pp}{\mathbf{P}}
\newcommand{\kk}{\mathbf{k}}

% 特殊な写像の記号
\newcommand{\ev}{\mathop{\mathrm{ev}}\nolimits} % 値写像
\newcommand{\pr}{\mathop{\mathrm{pr}}\nolimits} % 射影

% スクリプト体にするコマンド
%   例えば　{\mcal C} のように用いる
\newcommand{\mcal}{\mathcal}

% 花文字にするコマンド 
%   例えば　{\h C} のように用いる
\newcommand{\h}{\mathscr}

% ヒルベルト空間などの記号
\newcommand{\F}{\mcal{F}}
\newcommand{\X}{\mcal{X}}
\newcommand{\Y}{\mcal{Y}}
\newcommand{\Hil}{\mcal{H}}
\newcommand{\RKHS}{\Hil_{k}}
\newcommand{\Loss}{\mcal{L}_{D}}
\newcommand{\MLsp}{(\X, \Y, D, \Hil, \Loss)}

% 偏微分作用素の記号
\newcommand{\p}{\partial}

% 角カッコの記号 (内積は下にマクロがあります)
\newcommand{\lan}{\langle}
\newcommand{\ran}{\rangle}



% 圏の記号など
\newcommand{\Set}{{\bf Set}}
\newcommand{\Vect}{{\bf Vect}}
\newcommand{\FDVect}{{\bf FDVect}}
%\newcommand{\Ring}{{\bf Ring}}
\newcommand{\Ab}{{\bf Ab}}
\newcommand{\Mod}{\mathop{\mathrm{Mod}}\nolimits}
\newcommand{\Modf}{\mathop{\mathrm{Mod}^\mathrm{f}}\nolimits}
\newcommand{\CGA}{{\bf CGA}}
\newcommand{\GVect}{{\bf GVect}}
\newcommand{\Lie}{{\bf Lie}}
\newcommand{\dLie}{{\bf Liec}}



% 射の集合など
\newcommand{\Map}{\mathop{\mathrm{Map}}\nolimits} % 写像の集合
\newcommand{\Hom}{\mathop{\mathrm{Hom}}\nolimits} % 射集合
\newcommand{\End}{\mathop{\mathrm{End}}\nolimits} % 自己準同型の集合
\newcommand{\Aut}{\mathop{\mathrm{Aut}}\nolimits} % 自己同型の集合
\newcommand{\Mor}{\mathop{\mathrm{Mor}}\nolimits} % 射集合
\newcommand{\Ker}{\mathop{\mathrm{Ker}}\nolimits} % 核
\newcommand{\Img}{\mathop{\mathrm{Im}}\nolimits} % 像
\newcommand{\Cok}{\mathop{\mathrm{Coker}}\nolimits} % 余核
\newcommand{\Cim}{\mathop{\mathrm{Coim}}\nolimits} % 余像

% その他便利なコマンド
\newcommand{\dip}{\displaystyle} % 本文中で数式モード
\newcommand{\e}{\varepsilon} % イプシロン
\newcommand{\dl}{\delta} % デルタ
\newcommand{\pphi}{\varphi} % ファイ
\newcommand{\ti}{\tilde} % チルダ
\newcommand{\pal}{\parallel} % 平行
\newcommand{\op}{{\rm op}} % 双対を取る記号
\newcommand{\lcm}{\mathop{\mathrm{lcm}}\nolimits} % 最小公倍数の記号
\newcommand{\Probsp}{(\Omega, \F, \P)} 
\newcommand{\argmax}{\mathop{\rm arg~max}\limits}
\newcommand{\argmin}{\mathop{\rm arg~min}\limits}





% ================================
% コマンドを追加する場合のスペース 
%\newcommand{\OO}{\mcal{O}}



\makeatletter
\renewenvironment{proof}[1][\proofname]{\par
  \pushQED{\qed}%
  \normalfont \topsep6\p@\@plus6\p@\relax
  \trivlist
  \item[\hskip\labelsep
%        \itshape
         \bfseries
%    #1\@addpunct{.}]\ignorespaces
    {#1}]\ignorespaces
}{%
  \popQED\endtrivlist\@endpefalse
}
\makeatother

\renewcommand{\proofname}{Proof.}



%\renewcommand\proofname{\bf 証明} % 証明
\numberwithin{equation}{section}
\newcommand{\cTop}{\textsf{Top}}
%\newcommand{\cOpen}{\textsf{Open}}
\newcommand{\Op}{\mathop{\textsf{Op}}\nolimits}
\newcommand{\Ob}{\mathop{\textrm{Ob}}\nolimits}
\newcommand{\id}{\mathop{\mathrm{id}}\nolimits}
\newcommand{\pt}{\mathop{\mathrm{pt}}\nolimits}
\newcommand{\res}{\mathop{\rho}\nolimits}
\newcommand{\A}{\mcal{A}}
\newcommand{\B}{\mcal{B}}
\newcommand{\C}{\mcal{C}}
\newcommand{\D}{\mcal{D}}
\newcommand{\E}{\mcal{E}}
\newcommand{\G}{\mcal{G}}
%\newcommand{\H}{\mcal{H}}
\newcommand{\I}{\mcal{I}}
\newcommand{\J}{\mcal{J}}
\newcommand{\OO}{\mcal{O}}
\newcommand{\Ring}{\mathop{\textsf{Ring}}\nolimits}
\newcommand{\cAb}{\mathop{\textsf{Ab}}\nolimits}
%\newcommand{\Ker}{\mathop{\mathrm{Ker}}\nolimits}
\newcommand{\im}{\mathop{\mathrm{Im}}\nolimits}
\newcommand{\Coker}{\mathop{\mathrm{Coker}}\nolimits}
\newcommand{\Coim}{\mathop{\mathrm{Coim}}\nolimits}
\newcommand{\rank}{\mathop{\mathrm{rank}}\nolimits}
\newcommand{\Ht}{\mathop{\mathrm{Ht}}\nolimits}
\newcommand{\supp}{\mathop{\mathrm{supp}}\nolimits}
\newcommand{\colim}{\mathop{\mathrm{colim}}}
\newcommand{\Tor}{\mathop{\mathrm{Tor}}\nolimits}

\newcommand{\cat}{\mathscr{C}}

%筆記体
\newcommand{\cA}{\mcal{A}}
\newcommand{\cB}{\mcal{B}}
\newcommand{\cC}{\mcal{C}}
\newcommand{\cD}{\mcal{D}}
\newcommand{\cE}{\mcal{E}}
\newcommand{\cF}{\mcal{F}}
\newcommand{\cG}{\mcal{G}}
\newcommand{\cH}{\mcal{H}}
\newcommand{\cI}{\mcal{I}}
\newcommand{\cJ}{\mcal{J}}
\newcommand{\cK}{\mcal{K}}
\newcommand{\cL}{\mcal{L}}
\newcommand{\cM}{\mcal{M}}
\newcommand{\cN}{\mcal{N}}
\newcommand{\cO}{\mcal{O}}
\newcommand{\cP}{\mcal{P}}
\newcommand{\cQ}{\mcal{Q}}
\newcommand{\cR}{\mcal{R}}
\newcommand{\cS}{\mcal{S}}
\newcommand{\cT}{\mcal{T}}
\newcommand{\cU}{\mcal{U}}
\newcommand{\cV}{\mcal{V}}
\newcommand{\cW}{\mcal{W}}
\newcommand{\cX}{\mcal{X}}
\newcommand{\cY}{\mcal{Y}}
\newcommand{\cZ}{\mcal{Z}}


\newcommand{\scA}{\mathscr{A}}
\newcommand{\scB}{\mathscr{B}}
\newcommand{\scC}{\mathscr{C}}
\newcommand{\scD}{\mathscr{D}}
\newcommand{\scE}{\mathscr{E}}
\newcommand{\scF}{\mathscr{F}}
\newcommand{\scN}{\mathscr{N}}
\newcommand{\scO}{\mathscr{O}}
\newcommand{\scV}{\mathscr{V}}
\newcommand{\scU}{\mathscr{U}}


\newcommand{\ibA}{\mathop{\text{\textit{\textbf{A}}}}}
\newcommand{\ibB}{\mathop{\text{\textit{\textbf{B}}}}}
\newcommand{\ibC}{\mathop{\text{\textit{\textbf{C}}}}}
\newcommand{\ibD}{\mathop{\text{\textit{\textbf{D}}}}}
\newcommand{\ibE}{\mathop{\text{\textit{\textbf{E}}}}}
\newcommand{\ibF}{\mathop{\text{\textit{\textbf{F}}}}}
\newcommand{\ibG}{\mathop{\text{\textit{\textbf{G}}}}}
\newcommand{\ibH}{\mathop{\text{\textit{\textbf{H}}}}}
\newcommand{\ibI}{\mathop{\text{\textit{\textbf{I}}}}}
\newcommand{\ibJ}{\mathop{\text{\textit{\textbf{J}}}}}
\newcommand{\ibK}{\mathop{\text{\textit{\textbf{K}}}}}
\newcommand{\ibL}{\mathop{\text{\textit{\textbf{L}}}}}
\newcommand{\ibM}{\mathop{\text{\textit{\textbf{M}}}}}
\newcommand{\ibN}{\mathop{\text{\textit{\textbf{N}}}}}
\newcommand{\ibO}{\mathop{\text{\textit{\textbf{O}}}}}
\newcommand{\ibP}{\mathop{\text{\textit{\textbf{P}}}}}
\newcommand{\ibQ}{\mathop{\text{\textit{\textbf{Q}}}}}
\newcommand{\ibR}{\mathop{\text{\textit{\textbf{R}}}}}
\newcommand{\ibS}{\mathop{\text{\textit{\textbf{S}}}}}
\newcommand{\ibT}{\mathop{\text{\textit{\textbf{T}}}}}
\newcommand{\ibU}{\mathop{\text{\textit{\textbf{U}}}}}
\newcommand{\ibV}{\mathop{\text{\textit{\textbf{V}}}}}
\newcommand{\ibW}{\mathop{\text{\textit{\textbf{W}}}}}
\newcommand{\ibX}{\mathop{\text{\textit{\textbf{X}}}}}
\newcommand{\ibY}{\mathop{\text{\textit{\textbf{Y}}}}}
\newcommand{\ibZ}{\mathop{\text{\textit{\textbf{Z}}}}}

\newcommand{\ibx}{\mathop{\text{\textit{\textbf{x}}}}}

%\newcommand{\Comp}{\mathop{\mathrm{C}}\nolimits}
%\newcommand{\Komp}{\mathop{\mathrm{K}}\nolimits}
%\newcommand{\Domp}{\mathop{\mathsf{D}}\nolimits}%複体のホモトピー圏
%\newcommand{\Comp}{\mathrm{C}}
%\newcommand{\Komp}{\mathrm{K}}
%\newcommand{\Domp}{\mathsf{D}}%複体のホモトピー圏

\newcommand{\Comp}{\mathop{\mathrm{C}}\nolimits}
\newcommand{\Komp}{\mathop{\mathsf{K}}\nolimits}
\newcommand{\Domp}{\mathord{\mathsf{D}}}%\nolimits}
\newcommand{\Kompl}{\mathop{\mathsf{K}^\mathrm{+}}\nolimits}
\newcommand{\Kompu}{\mathop{\mathsf{K}^\mathrm{-}}\nolimits}
\newcommand{\Kompb}{\mathop{\mathsf{K}^\mathrm{b}}\nolimits}
\newcommand{\Dompl}{\mathop{\mathsf{D}^\mathrm{+}}\nolimits}
\newcommand{\Dompu}{\mathop{\mathsf{D}^\mathrm{-}}\nolimits}
\newcommand{\Dompb}{\mathop{\mathsf{D}^\mathrm{b}}\nolimits}
\newcommand{\Dompbf}{\mathop{\mathsf{D}_\mathrm{f}^\mathrm{b}}\nolimits}
\newcommand{\Domplb}{\mathop{\mathsf{D}^\mathrm{lb}}\nolimits}




\newcommand{\CCat}{\Comp(\cat)}
\newcommand{\KCat}{\Komp(\cat)}
\newcommand{\DCat}{\Domp(\cat)}%圏Cの複体のホモトピー圏
\newcommand{\HOM}{\mathop{\mathscr{H}\hspace{-2pt}om}\nolimits}%内部Hom
\newcommand{\RHOM}{\mathop{\mathrm{R}\hspace{-1.5pt}\HOM}\nolimits}

%\newcommand{\muS}{\mathop{\mathrm{SS}}\nolimits}
\newcommand{\muS}[1][]{\mathop{\mathrm{SS}}\nolimits_{#1}}
\newcommand{\RG}{\mathop{\mathrm{R}\hspace{-0pt}\Gamma}\nolimits}
\newcommand{\RHom}{\mathop{\mathrm{R}\hspace{-1.5pt}\Hom}\nolimits}
%\newcommand{\Rder}{\mathrm{R}}
\newcommand{\Rder}[1][]{\mathop{\mathrm{R}^{#1}\hspace{-2pt}}\nolimits}

\newcommand{\simar}{\mathrel{\overset{\sim}{\rightarrow}}}%同型右矢印
\newcommand{\simarr}{\mathrel{\overset{\sim}{\longrightarrow}}}%同型右矢印
\newcommand{\simra}{\mathrel{\overset{\sim}{\leftarrow}}}%同型左矢印
\newcommand{\simrra}{\mathrel{\overset{\sim}{\longleftarrow}}}%同型左矢印

\newcommand{\hocolim}{{\mathrm{hocolim}}}
\newcommand{\indlim}[1][]{\mathop{\varinjlim}\limits_{#1}}
\newcommand{\sindlim}[1][]{\smash{\mathop{\varinjlim}\limits_{#1}}\,}
\newcommand{\Pro}{\mathrm{Pro}}
\newcommand{\Ind}{\mathrm{Ind}}
\newcommand{\prolim}[1][]{\mathop{\varprojlim}\limits_{#1}}
\newcommand{\sprolim}[1][]{\smash{\mathop{\varprojlim}\limits_{#1}}\,}

\newcommand{\Sh}{\mathrm{Sh}}
\newcommand{\PSh}{\mathrm{PSh}}

\newcommand{\rmD}{\mathrm{D}}

\newcommand{\Lloc}[1][]{\mathord{\mathcal{L}^1_{\mathrm{loc},{#1}}}}
\newcommand{\ori}{\mathord{\mathrm{or}}}
\newcommand{\Db}{\mathord{\mathscr{D}b}}

\newcommand{\codim}{\mathop{\mathrm{codim}}\nolimits}



\newcommand{\gld}{\mathop{\mathrm{gld}}\nolimits}
\newcommand{\wgld}{\mathop{\mathrm{wgld}}\nolimits}


\newcommand{\tens}[1][]{\mathbin{\otimes_{\raise1.5ex\hbox to-.1em{}{#1}}}}
\newcommand{\etens}{\mathbin{\boxtimes}}
\newcommand{\ltens}[1][]{\mathbin{\overset{\mathrm{L}}\tens}_{#1}}
\newcommand{\mtens}[1][]{\mathbin{\overset{\mathrm{\mu}}\tens}_{#1}}
\newcommand{\lltens}[1][]{{\mathop{\tens}\limits^{\mathrm{L}}_{#1}}}
\newcommand{\letens}{\overset{\mathrm{L}}{\etens}}
\newcommand{\detens}{\underline{\etens}}
\newcommand{\ldetens}{\overset{\mathrm{L}}{\underline{\etens}}}
\newcommand{\dtens}[1][]{{\overset{\mathrm{L}}{\underline{\otimes}}}_{#1}}

\newcommand{\blk}{\mathord{\ \cdot\ }}
\newcommand{\mres}[2][]{{\left.{#1}\right\rvert}_{#2}}


%\newcommand{\hocolim}{{\mathrm{hocolim}}}
%\newcommand{\indlim}[1][]{\mathop{\varinjlim}\limits_{#1}}
%\newcommand{\sindlim}[1][]{\smash{\mathop{\varinjlim}\limits_{#1}}\,}
%\newcommand{\Pro}{\mathrm{Pro}}
%\newcommand{\Ind}{\mathrm{Ind}}
%\newcommand{\prolim}[1][]{\mathop{\varprojlim}\limits_{#1}}
%\newcommand{\sprolim}[1][]{\smash{\mathop{\varprojlim}\limits_{#1}}\,}
\newcommand{\proolim}[1][]{\mathop{\text{\rm``{$\varprojlim$}''}}\limits_{#1}}
\newcommand{\sproolim}[1][]{\smash{\mathop{\rm``{\varprojlim}''}\limits_{#1}}}
\newcommand{\inddlim}[1][]{\mathop{\text{\rm``{$\varinjlim$}''}}\limits_{#1}}
\newcommand{\sinddlim}[1][]{\smash{\mathop{\text{\rm``{$\varinjlim$}''}}\limits_{#1}}\,}
\newcommand{\ooplus}{\mathop{\text{\rm``{$\oplus$}''}}\limits}
\newcommand{\bbigsqcup}{\mathop{``\bigsqcup''}\limits}
\newcommand{\bsqcup}{\mathop{``\sqcup''}\limits}
\newcommand{\dsum}[1][]{\mathbin{\oplus_{#1}}}

\newcommand{\Fct}{\mathop{\mathsf{Fct}}\nolimits}

\newcommand{\pb}[1][]{\mathbin{\mathop{\times}\limits_{#1}}}
\newcommand{\dpb}[1][]{\mathbin{\mathop{\times}\limits_{#1}}}
\newcommand{\tpb}[1][]{\mathbin{\mathop{\times}\nolimits_{#1}}}





%================================================
% 自前の定理環境
%   https://mathlandscape.com/latex-amsthm/
% を参考にした
\newtheoremstyle{mystyle}%   % スタイル名
    {5pt}%                   % 上部スペース
    {5pt}%                   % 下部スペース
    {}%              % 本文フォント
    {}%                  % 1行目のインデント量
    {\bfseries}%                      % 見出しフォント
    {.}%                     % 見出し後の句読点
    {12pt}%                     % 見出し後のスペース
    {\thmname{#1}\thmnumber{ #2}\thmnote{{\hspace{2pt}\normalfont (#3)}}}% % 見出しの書式

\theoremstyle{mystyle}
\newtheorem{AXM}{公理}%[section]
\newtheorem{DFN}[AXM]{定義}
\newtheorem{THM}[AXM]{定理}
\newtheorem*{THM*}{定理}
\newtheorem{PRP}[AXM]{命題}
\newtheorem{LMM}[AXM]{補題}
\newtheorem{CRL}[AXM]{系}
\newtheorem{EG}[AXM]{例}
\newtheorem*{EG*}{例}
\newtheorem{RMK}[AXM]{注意}
\newtheorem{CNV}[AXM]{約束}
\newtheorem{CMT}[AXM]{コメント}
\newtheorem*{CMT*}{コメント}
\newtheorem{NTN}[AXM]{記号}

% 定理環境ここまで
%====================================================

\usepackage{framed}
\definecolor{lightgray}{rgb}{0.75,0.75,0.75}
\renewenvironment{leftbar}{%
  \def\FrameCommand{\textcolor{lightgray}{\vrule width 4pt} \hspace{10pt}}% 
  \MakeFramed {\advance\hsize-\width \FrameRestore}}%
{\endMakeFramed}
\newenvironment{redleftbar}{%
  \def\FrameCommand{\textcolor{lightgray}{\vrule width 1pt} \hspace{10pt}}% 
  \MakeFramed {\advance\hsize-\width \FrameRestore}}%
 {\endMakeFramed}



\renewcommand{\Re}{\mathop{\mathrm{Re}}\nolimits}

% =================================





% ---------------------------
% new definition macro
% ---------------------------
% 便利なマクロ集です

% 内積のマクロ
%   例えば \inner<\pphi | \psi> のように用いる
\def\inner<#1>{\langle #1 \rangle}

% ================================
% マクロを追加する場合のスペース 

%=================================



% 位相空間
\newcommand{\Int}[1][]{\mathop{\mathrm{Int}}\nolimits_{#1}}
\newcommand{\Cl}[1][]{\mathop{\mathrm{Cl}}\nolimits_{#1}}
\newcommand{\Bd}[1][]{\mathop{\mathrm{Bd}}\nolimits_{#1}}
\newcommand{\bd}{\mathord{\partial}}
\newcommand{\Fr}[1][]{\mathop{\mathrm{Fr}}\nolimits_{#1}}


\newcommand{\mtan}[1][]{T\hspace{-1.75pt}{#1}}
\newcommand{\mtp}[1][]{{}^{t\!}{#1}} % transpose of the morphism
%\newcommand{\mpb}[1][]{{}^{t\!}{#1}'} % pull-back of the morphism
\newcommand{\mpb}[1][]{{#1}_{\pi}} % pull-back of the morphism
\newcommand{\mdiff}[1][]{{#1}_{d}} % pull-back of the morphism
\newcommand{\mptan}[1][]{{#1}_{\tau}} % pull-back of the morphism


%\newcommand{\pb}[1][]{\mathbin{\mathop{\times}\limits_{#1}}}

\newcommand{\cotb}[1][]{T^{\ast}{#1}}
\newcommand{\conm}[2][]{T_{#1}^{\hspace{0.7pt}\ast}{#2}}
\newcommand{\norb}[2][]{T_{#1}{#2}}

\newcommand{\cotz}[1][]{T^{\ast}_{#1}{#1}}

\newcommand{\cotn}[1][]{\dot{T}^{\ast}{#1}}


\newcommand{\muhom}{\mu hom}
\newcommand{\muHom}[1][]{\mathrm{Hom}^\mu_{\raise1.5ex\hbox to.1em{}#1}}

\newcommand{\mucat}[2][]{\mathsf{D}^{\mathrm{b}}_{#1}(#2)}
%\newcommand{\mucat}[2][]{\mathsf{D}^{\mathrm{b}}_{#1}(#2)}

\newcommand{\conv}[1][]{\mathbin{\mathop{\circ}\limits_{#1}}}


%\newcommand{\lopen}{\mathopen{]}}
%\newcommand{\lclosed}{\mathopen{[}}
%\newcommand{\ropen}{\mathclose{[}}
%\newcommand{\rclosed}{\mathopen{]}}
\newcommand{\lopen}{\mathopen{(}}
\newcommand{\lclosed}{\mathopen{[}}
\newcommand{\ropen}{\mathclose{)}}
\newcommand{\rclosed}{\mathopen{]}}

\newcommand{\sm}{\raisebox{2.33pt}{~\rule{6.4pt}{1.3pt}~}}
%\newcommand{\smi}{\mathbin{\raisebox{2.33pt}{\rule{6.4pt}{1.3pt}}}}
\newcommand{\smi}{-}

\newcommand{\smid}{\mathrel{;}}

%\figre名を英語にする
\renewcommand{\figurename}{Fig.}
%\参考文献名を英語にする
\renewcommand{\refname}{References}



%\theoremstyle{definition}
\newtheorem{defn}{定義}
\newtheorem{thm}[defn]{定理}
\newtheorem{prop}[defn]{命題}
\newtheorem{lemma}[defn]{補題}
\newtheorem{rem}[defn]{注意}

%\documentclass[a4paper]{jarticle}
% として jarticle を使う場合は，上の \usepackage の部分を
%\usepackage[jslayout]{syuron2}
% のように [jslayout] を追加する．

%\documentclass[a4paper]{jarticle}
% として jarticle を使う場合は，上の \usepackage の部分を
%\usepackage[jslayout]{gaiyo2}
% のように [jslayout] を追加する．

% 以下は修論のフォーマットを参照して記入する．
\title{Microlocal Morse Theory and Truncated Microsupport\\（超局所モース理論と打切り超局所台）}
\author{大柴寿浩}
\date{2026.3} % 例. 2024年3月卒業の場合は「2024.3」を記入する．
\kei{解析系} % 自分の系を記入する．
\shisho{本多尚文 教授} % 教授・准教授・助教の別に注意

\begin{document}
\begin{gaiyo}
  %指導教員の指導のもとに修士論文の概要を作成する．概要の形式と
  %内容は指導教員とよく相談するべきである．
  %修士論文概要は，修士論文
  %審査会に際して各系の教員が参照しながら発表を聞く．
  %従って，各系の中である程度広い範囲で理解されるように書く
  %必要がある．
  %ここでは参考のために一般的な形式を紹介する．
  %
  %まず修士論文の基本的な内容を短く紹介する．これは
  %系のメンバーが広く理解できるように書く．次に，修士論文の内容に関する背景を
  %簡潔に紹介する．これは関連分野に属していれば理解
  %を述べる．
  %修士論文に独自の結果を含むか否か，総合報告であるかどうかについて，
  %ここまでに明示しておく．
  %
  %以上は主要な結果を述べるための準備である．以下，
  %主要な結果を要約する．主定理を正確に記述し，必要な定義などを記述する．
  %次に，修士論文に従って，主要な結果の証明の方針などを記述する．
  %%必要であれば，最後に参考文献を挙げる．
  %

  %本修士論文の目的は，
  超局所層理論は，超局所解析の層理論における類似である．
  まず，超局所解析について復習し，「超局所」という言葉の意味を説明する．

  %古典的な微分方程式の問題を考えよう．
  \(n\)次元ユークリッド空間\(\mathbf{R}^{n}\)内の
  開集合\(\varOmega\)を考える．
  \(\varOmega\)の上に余接束\(
    T^{\ast}\varOmega=\varOmega\times \mathbf{R}^{n}
  \)という空間が考えられる．これは物理学における相空間の数学的なモデルである．
  簡単のためにここでは第2成分から原点を除いた\(
    \dot{T}^{\ast}\varOmega=\varOmega\times (\mathbf{R}^{n}-\{0\})
  \)を考えることにする．
  このとき，自然な射影\(T^{\ast}\varOmega\to \varOmega\)や
  その制限\(\dot{T}^{\ast}\varOmega\to \varOmega\)を\(\pi\)で表す．
  \(\varOmega\)上の
  微分方程式\(Pu=f\)を考えよう．
  ここで\(u\)は未知関数，\(f\)は既知関数，\(P\)は実解析関数を係数とする
  線形偏微分作用素とする．
  一般に佐藤超関数 (hyperfunction)
  %（これはシュワルツのdistributionを
  %含む超関数概念である）
  \(u\)に対し，\(u\)の特異台 (singular support) が
  \(
    \mathrm{singsupp}(u)
    =\left\{x\in\varOmega;\ 
    \text{\(u(x)\)は\(x\)で解析的でない}
    \right\}
  \)として定義される．これは定空間の部分集合である．一方，
  \(u\)の特異性スペクトル (singularity spectrum), 
  あるいは（解析的）波面集合 ((analytic) wavefront set) と呼ばれる
  余接束の部分集合\(\mathrm{S.S.}(u)\)が（正確な定義は省くが）\(
    \pi(\mathrm{S.S.}(u))=\mathrm{singsupp}(u)
  \)を満たすように定まる．
  さらに，微分作用素\(P\)に対しても特性多様体 (characteristic variety) と
  呼ばれる余接束の部分集合\(\mathrm{Char}(P)\)が\(
    \mathrm{Char}(P)=\{
      (x;\xi)\in\dot{T}^{\ast}\varOmega;\ 
      \sigma(P)(x;\xi)=0
    \}
  \)として定義される．ここで\(\sigma(P)(x;\xi)\)は，
  \(P\)の微分の部分を不定元に置き換えた解析関数係数の多項式である．
  これを\(P\)の主要表象と呼ぶ．
  以上の準備のもと，佐藤の基本定理と呼ばれる次の評価が知られている：\[
    \mathrm{S.S.}(u)\subset \mathrm{Char}(P)\cup\mathrm{S.S.}(u).
  \]これを\(P\)が楕円型作用素，
  すなわち\(\mathrm{Char}(P)=\varnothing\)となる
  作用素と\(f=0\)すなわち斉次方程式に対して適用する．
  すると\(\pi\)による射影を考えれば
  \(\mathrm{singsupp}(u)\subset\varnothing\)を得る．
  すなわち，楕円斉次方程式の佐藤超関数解は実解析的である．
  これは楕円型斉次方程式の解の内部正則性を意味する．
  このように，考えたい数学的対象に対し，
  余接束の部分集合をなんらかの方法で定義し，
  それに関する評価を得ることで，その対象の特異性を調べることができる．
  「超局所」とは，このように
  空間上の対象物を余接束上にもちあげて考えることをいう．
  超局所解析の基礎理論は佐藤幹夫とその周辺人物により，1970年台に確立された．

  1980年台になると，その層理論における類似が
  柏原--Schapira により開発された．
  その中心概念の一つに，層の超局所台 (microsupport) がある．
  多様体上の層\(F\)に対し，
  \(F\)の超局所台\(\mathrm{SS}(F)\)が
  余接束の部分集合として定義される．超局所台を用いると，
  \(F\)のコホモロジー\(H^{j}(F)\)の伝播が調べられる．
  特に層を入力とする関手が同型になる条件が超局所台を用いて述べることができる．
  例えば，多様体の射\(f\colon X\to Y\)に対し，
  逆像と双対化複体とのテンソル積\(
    f^{-1}(\cdot)\otimes \omega_{X/Y}
  \)とねじれ逆像\(f^{!}(\cdot)\)との間には自然な射\(
    f^{-1}(\cdot)\otimes \omega_{X/Y}
    \to f^{!}(\cdot)
  \)が存在する．これが同型になるためには，
  \(f\)が\(F\)に関して非特性的であればよいことが知られている．
  この「非特性的」という条件は，\(f\)が閉うめこみのときには，
  \(X\)の\(Y\)に関する余法束と\(\mathrm{SS}(F)\)との共通部分が
  零切断に含まれるということになる．
  このような理論を体系化したのが超局所層理論である．

  超局所層理論は\(\mathcal{D}\)加群へも応用される．
  \(\mathcal{D}\)加群は線形微分方程式系の一般化である．
  複素多様体\(X\)を考える．\(X\)上の
  \(\mathcal{D}_{X}\)加群\(\mathcal{M}\)に対しても
  特性多様体\(\mathrm{Char}(\mathcal{M})\)が余接束の部分集合として定まる．
  このとき，\(\mathcal{M}\)の
  正則関数解の層\(
    \mathcal{S}ol(\mathcal{M})\coloneqq 
    \mathrm{R}\mathcal{H}om_{\mathcal{D}_{X}}(\mathcal{M},\mathcal{O}_{X})
  \)を考えると，
  \(\mathrm{Char}(\mathcal{M})=\mathrm{SS}(\mathcal{S}ol(\mathcal{M}))\)が成り立つ．
  さて，上で紹介した\(
    f^{-1}(\cdot)\otimes \omega_{X/Y}
    \to f^{!}(\cdot)
  \)が同型になるための条件を実解析多様体\(M\)から
  その複素近傍\(X\)への閉うめこみに\(i\colon M\hookrightarrow X\)対して適用しよう．
  しかも\(\mathcal{M}\)が楕円型，
  すなわち\(\mathrm{Char}(\mathcal{M})\)と
  零切断を除いた余法束との共通部分が空であるとする．
  このとき，\(i\)が\(\mathcal{S}ol(\mathcal{M})\)に関して非特性的ならば，
  \[
    i^{-1}\mathcal{S}ol(\mathcal{M})\otimes \omega_{M/X}
    \overset{\sim}{\to}
    i^{!}\mathcal{S}ol(\mathcal{M})
  \]である．これから，
  実解析関数の解の層と佐藤超関数の解の層が同型であることがわかる．
  すなわち楕円型正則性定理の（超局所）層理論的な証明が得られる．

  以上のように，超局所台を用いることで解析的な応用も持つ種々の主張が示される．
  ところが，境界値問題のような場合には，層のすべての次数の
  コホモロジーが伝播せず，
  特定の次数以下のコホモロジーしか伝播しないという現象が知られている．
  このような場合を扱うために，柏原--Monteiro-Fernandes--Schapira は
  2003年に，超局所台を精密化した打切り超局所台 (truncated microsupport) を導入した．
  打切り超局所台を用いることで，次数の条件がついた伝播現象を取り扱うことが
  できるようになった．

  本修士論文では，この打切り超局所台を用いて，層係数Morse理論を研究する．
  本論文では2つの主定理を示す．一つは非特性変形補題と呼ばれる，超局所層理論
  における基本定理の打切り版である．
  \begin{THM}
    \(X\)をハウスドルフ空間で，任意の開集合がパラコンパクトであるものとする．
    \(F\in\Dompl(\kk_X)\)とし，
    \(\left(U_{t}\right)_{t\in\rr}\)を\(X\)の開集合の族とする．
    \(k\in\zz\)とする．
    \(s\in\rr\)に対し\(
      Z_{s}\coloneqq \bigcap_{t>s}\overline{U_t\smi U_s}
    \)とおく．
    次を仮定する．
    \begin{enumerate}[(i)]
      \item \label{item:t-nonchar1}
      任意の\(t\in\rr\)に対し\(U_{t}=\bigcup_{s<t}U_{s}\)である．
      \item \label{item:t-nonchar2}
      任意の実数\(s\leqq t\)に対し\(
        \overline{U_{t}\smi U_{s}}\cap \supp{F}
      \)はコンパクトである．
      \item \label{item:t-nonchar3}
      任意の実数\(s\in\rr\)と任意の点\(x\in Z_{s}\smi U_{s}\)に対し，
      \(j\leqq k\)ならば次が成り立つ．
      \[
        (H^{j}_{X\smi U_{s}}(F))_{x}=0. 
      \]
      \item \label{item:t-nonchar4} 
      任意の実数\(s< t\)に対し，\(Z_s\subset U_{t}\)が成り立つ．
    \end{enumerate}
    このとき，任意の\(t\in\rr\)に対し，\(j<k\)ならば，制限射\[
      H^{j}\left(\bigcup_{s\in\rr}U_{s};F\right)
      \simar 
      H^{j}\left(U_{t};F\right)
    \]は同型であり，\(j=k\)ならば\[
      H^{k}\left(\bigcup_{s\in\rr}U_{s};F\right)
      \hookrightarrow 
      H^{k}\left(U_{t};F\right)
    \]は単射である．
  \end{THM}
  もう一つの主定理は，その打切り版非特性変形補題のMorse理論への応用である．
  \begin{THM}\label{thm:t-morse}
    \(F\in\Dompb(\kk_M)\)とし，
    \(\psi\colon M\to \rr\)を実数値\(C^1\)関数で
    \(\supp(F)\)上固有なものとする．
    \(a<b\)を実数とし\(k\in\zz\)とする．
    \(
      a\leqq \psi(x)<b
    \)に対し，\(d\psi(x)\notin \muS[k](F)\)が成り立つとする．
    実数\(t\in \rr\)に対し\(
      M_t\coloneqq\psi^{-1}\left((-\infty,t)\right)
    \)とおく．このとき制限射
    \[
      H^{j}(M_{b};F)\to H^{j}(M_{a};F)
    \]
    は\(j< k\)ならば同型であり\(j=k\)ならば単射である．
  \end{THM}
  この2つが本論文での独自の結果である．

  本論文の以下のとおりである．
  第1節は導入である．ここでは理論の背景と論文の概要を述べてある．
  第2節では，導来圏の理論をもとに
  層理論とホモロジー代数における基本的な事実の復習を行う．
  特に，Mittag-Leffler条件を用いた射影極限に関する主張と，
  開集合とは限らない集合上の
  層係数コホモロジーの帰納極限による近似は主定理の証明で重要な役割を演ずる．
  第3節では打切り超局所台の定義とその順像公式を紹介する．
  第4節では，第1の主定理である打切り版非特性変形補題の主張を述べる．
  打切り版の主張は技術的な問題から元来のものとは少し状況設定が異なる．
  そのため，元来の主張は満たすが，我々の場合の仮定を満たさない例を1つ挙げ，
  我々の場合の仮定を満たす例を3つ挙げた．これらは図も併記した．
  これらの例を述べたのち，主定理の証明を行う．定理の証明中で鍵となる補題が
  補題4.4である．この補題の証明は長く技術的なので，第6節に掲載した．
  第5節では，第2の主定理である打切り版超局所モースの補題を証明する．
  はじめに打切り版非特性変形補題を1次元ユークリッド空間の開区間とその上の層
  に適用し，一般の多様体上の主張を打切り超局所台の順像に関する関手的性質を
  用いて1次元ユークリッド空間の場合に帰着させる．
\end{gaiyo}

超局所層理論は，超局所解析の層理論における類似である．
まず，超局所解析について復習し，「超局所」という言葉の意味を説明する．
\(\rr^{n}\)の開集合\(\varOmega\)を考える．
\(\varOmega\)の上に余接束\(
  T^{\ast}\varOmega=\varOmega\times \mathbf{R}^{n}
\)という空間が考えられる．
簡単のためにここでは\(
  \dot{T}^{\ast}\varOmega=\varOmega\times (\mathbf{R}^{n}-\{0\})
\)を考えることにする．
\(\varOmega\)上の
微分方程式\(Pu=f\)を考えよう．
ここで\(u\)は未知関数，\(f\)は既知関数，\(P\)は実解析係数
線形偏微分作用素とする．
一般に佐藤超関数
\(u\)に対し，\(u\)の特異台が
\(
  \mathrm{singsupp}(u)
  =\left\{x\in\varOmega;\ 
  \text{\(u(x)\)は\(x\)で解析的でない}
  \right\}
\)として定義される．一方，
\(u\)の特異性スペクトル，
あるいは（解析的）波面集合と呼ばれる
余接束の部分集合\(\mathrm{S.S.}(u)\)が（正確な定義は省くが）\(
  \pi(\mathrm{S.S.}(u))=\mathrm{singsupp}(u)
\)を満たすように定まる．
微分作用素\(P\)に対しても特性多様体 (characteristic variety) と
呼ばれる余接束の部分集合\(\mathrm{Char}(P)\)が\(
  \mathrm{Char}(P)=\{
    (x;\xi)\in\dot{T}^{\ast}\varOmega;\ 
    \sigma(P)(x;\xi)=0
  \}
\)として定義される．ここで\(\sigma(P)(x;\xi)\)は，
\(P\)の主要表象である．
このとき，佐藤の基本定理：\(
  \mathrm{S.S.}(u)\subset \mathrm{Char}(P)\cup\mathrm{S.S.}(u).
\)が成り立つ．これを\(P\)が楕円型作用素
と\(f=0\)に対して適用する．
すると\(\pi\)による射影を考えれば
\(\mathrm{singsupp}(u)\subset\varnothing\)を得る．
すなわち，楕円斉次方程式の佐藤超関数解は実解析的である．
「超局所」とは，このように
空間上の対象物を余接束上にもちあげて考えることをいう．
超局所解析の基礎理論は佐藤幹夫とその周辺人物により，1970年台に確立された．

1980年台になると，その層理論における類似が
柏原--Schapira により開発された\cite{KS90}．
その中心概念の一つに，層の超局所台がある．
多様体上の層\(F\)に対し，
\(F\)の超局所台\(\mathrm{SS}(F)\)が
余接束の部分集合として定義される．超局所台を用いると，
\(F\)のコホモロジー\(H^{j}(F)\)の伝播が調べられる．
特に層を入力とする関手が同型になる条件が超局所台を用いて述べることができる．
例えば，多様体の射\(f\colon X\to Y\)に対し，
自然な射\(
  f^{-1}(\cdot)\otimes \omega_{X/Y}
  \to f^{!}(\cdot)
\)が存在する．これが同型になるためには，
\(f\)が\(F\)に関して非特性的であればよい．
これは\(f\)が閉うめこみのときには，
\(X\)の\(Y\)に沿った余法束と\(\mathrm{SS}(F)\)との共通部分が
零切断に含まれるということである．
超局所層理論は\(\mathcal{D}\)加群へも応用される．
\(\mathcal{D}\)加群は線形微分方程式系の一般化である．
複素多様体\(X\)上の
\(\mathcal{D}_{X}\)加群\(\mathcal{M}\)に対しても
特性多様体\(\mathrm{Char}(\mathcal{M})\)が定まる．
このとき，\(\mathcal{M}\)の
正則関数解の層\(
  \mathcal{S}ol(\mathcal{M})\coloneqq 
  \mathrm{R}\mathcal{H}om_{\mathcal{D}_{X}}(\mathcal{M},\mathcal{O}_{X})
\)を考えると，
\(\mathrm{Char}(\mathcal{M})=\mathrm{SS}(\mathcal{S}ol(\mathcal{M}))\)が成り立つ．
さて，上で紹介した\(
  f^{-1}(\cdot)\otimes \omega_{X/Y}
  \simar f^{!}(\cdot)
\)を実解析多様体\(M\)から
その複素近傍\(X\)への閉うめこみ\(i\colon M\hookrightarrow X\)と
楕円型\(\cD_{X}\)加群\(\mathcal{M}\)，
すなわち\(\mathrm{Char}(\mathcal{M})\)と
零切断を除いた余法束との共通部分が空なものに適用することで次を得る．
\[
  i^{-1}\mathcal{S}ol(\mathcal{M})\otimes \omega_{M/X}
  \overset{\sim}{\to}
  i^{!}\mathcal{S}ol(\mathcal{M})
\]これから，
実解析関数の解の層と佐藤超関数の解の層が同型であることがわかる．

以上のように，超局所台を用いることで解析的にも有用な種々の主張が示される．
ところが，境界値問題のような場合には，層のすべての次数の
コホモロジーが伝播せず，
特定の次数以下のコホモロジーしか伝播しないという現象が知られている．
このような場合を扱うために，柏原--Monteiro-Fernandes--Schapira は
2003年に，超局所台を精密化した打切り超局所台を導入した\cite{KFS03a}．
打切り超局所台を用いることで，次数の条件がついた伝播現象を取り扱うことが
できるようになった．

本修士論文の目的は，この打切り超局所台を用いて層係数Morse理論を研究することである．
特に以下の2つの主定理を示す．一つは非特性変形補題と呼ばれる，超局所層理論
における基本定理の打切り版である．
\begin{THM}
  \(X\)をハウスドルフ空間で，任意の開集合がパラコンパクトであるものとする．
  \(F\in\Dompl(\kk_X)\)とし，
  \(\left(U_{t}\right)_{t\in\rr}\)を\(X\)の開集合の族とする．
  \(k\in\zz\)とする．
  \(s\in\rr\)に対し\(
    Z_{s}\coloneqq \bigcap_{t>s}\overline{U_t\smi U_s}
  \)とおく．
  次を仮定する．
  \begin{enumerate}[(i)]
    \item \label{item:t-nonchar1}
    任意の\(t\in\rr\)に対し\(U_{t}=\bigcup_{s<t}U_{s}\)である．
    \item \label{item:t-nonchar2}
    任意の実数\(s\leqq t\)に対し\(
      \overline{U_{t}\smi U_{s}}\cap \supp{F}
    \)はコンパクトである．
    \item \label{item:t-nonchar3}
    任意の実数\(s\in\rr\)と任意の点\(x\in Z_{s}\smi U_{s}\)に対し，
    \(j\leqq k\)ならば次が成り立つ．
    \[
      (H^{j}_{X\smi U_{s}}(F))_{x}=0. 
    \]
    \item \label{item:t-nonchar4} 
    任意の実数\(s< t\)に対し，\(Z_s\subset U_{t}\)が成り立つ．
  \end{enumerate}
  このとき，任意の\(t\in\rr\)に対し，\(j<k\)ならば，制限射\[
    H^{j}\left(\bigcup_{s\in\rr}U_{s};F\right)
    \simar 
    H^{j}\left(U_{t};F\right)
  \]は同型であり，\(j=k\)ならば\[
    H^{k}\left(\bigcup_{s\in\rr}U_{s};F\right)
    \hookrightarrow 
    H^{k}\left(U_{t};F\right)
  \]は単射である．
\end{THM}
もう一つの主定理は，その打切り版非特性変形補題のMorse理論への応用である．
\begin{THM}\label{thm:t-morse}
  \(F\in\Dompb(\kk_M)\)とし，
  \(\psi\colon M\to \rr\)を実数値\(C^1\)関数で
  \(\supp(F)\)上固有なものとする．
  \(a<b\)を実数とし\(k\in\zz\)とする．
  \(
    a\leqq \psi(x)<b
  \)に対し，\(d\psi(x)\notin \muS[k](F)\)が成り立つとする．
  実数\(t\in \rr\)に対し\(
    M_t\coloneqq\psi^{-1}\left((-\infty,t)\right)
  \)とおく．このとき制限射
  \[
    H^{j}(M_{b};F)\to H^{j}(M_{a};F)
  \]
  は\(j< k\)ならば同型であり\(j=k\)ならば単射である．
\end{THM}
この2つが本論文での独自の結果である．

本論文の構成は以下のとおりである．
第1節は導入である．ここでは理論の背景と論文の概要を述べる．
第2節では，導来圏を用いた
層係数こホモロジーの復習を行う．
特に，Mittag-Leffler条件に関する主張と，
層係数コホモロジーの帰納極限による近似が重要である．
第3節では打切り超局所台の定義とその順像公式を紹介する．
第4節では，第1の主定理である打切り版非特性変形補題の主張を述べる．
打切り版の主張は元のものとは状況設定が異なる．
そのため，元の主張は満たすが，打切り版の仮定を満たさない例を1つ，
満たす例を3つ，図とともに挙げる．
これらの例を述べたのち，主定理の証明を行う．定理の証明中で鍵となる補題が
補題4.4である．この補題の証明は長く技術的なので，第6節に掲載した．
第5節では，第2の主定理である打切り版超局所モースの補題を証明する．
はじめに打切り版非特性変形補題を1次元ユークリッド空間の開区間とその上の層
に適用し，一般の多様体上の主張を打切り超局所台の順像に関する関手的性質を
用いて1次元ユークリッド空間の場合に帰着させる．
%===============================================
% 参考文献スペース
%===============================================
\begin{thebibliography}{20} 
  \bibitem[KFS03a]{KFS03a} Masaki Kashiwara, Teresa Monteiro Fernandes, Pierre Schapira, 
    \textit{Truncated Microsupport and Holomorphic Solutions 
    of \(D\)-modules}, 
    Ann. Scient. \'Ec. Norm. Sup., s\'erie 4, tome 36, 
    pp.583--399, 2003.
  \bibitem[KS90]{KS90} Masaki Kashiwara, Pierre Schapira, 
    \textit{Sheaves on Manifolds}, 
    Grundlehren der Mathematischen Wissenschaften, 292, Springer, 1990.
\end{thebibliography}
%===============================================

\end{document}
